\section{总结}

\subsection{可扩展之处}

\begin{enumerate}
  \item \textbf{扩展连续视觉对准}：将现有板端连续对准扩展到小螺丝，结合运动中的图像更新，缩短接近与对准之间的等待。
  \item \textbf{扩展零件与工位}：增加垫片、销钉等零件的识别样本和抓取参数，将示教记录接入槽位配置，简化工位调整。
  \item \textbf{进一步利用 PL}：在现有掩码核之后评估轮廓和矩计算的硬件实现，并优化位置读取与 CAN 周期调度的配合。
  \item \textbf{复用工程工具}：将 CAN 时序分析、NPU 部署和工位运维流程整理为通用工具，支持其他 Kria 与 CAN 执行器项目。
\end{enumerate}

\subsection{心得体会}

本项目首先解决的是计算任务的分工。模型推理和轨迹规划适合放在 AI PC，需要接近执行端的腕部定位放在 KV260 的 ARM，周期执行和设定检查交给 PL。最初由 ARM 软件逐拍发送 CAN 设定时，控制服务暂停 20\ms 会使最长命令间隔达到 21.435\ms；改为 PL 执行后，同一对照中的最长间隔为 5.014\ms。随后用逻辑分析仪测量总线，确认硬件周期在真实通信中得到保持。

NPU 部署的价值体现在并发运行。单独检测时 GPU 更快，但大模型同时生成时，GPU 检测尾延迟明显增加。将持续检测放到 NPU 后，系统可以在解析指令、复审结果的同时保持视觉更新。量化中遇到过类别分数因统一 INT8 尺度而失真的问题，最终保留解码节点的浮点计算，再逐点比较两种后端的抓取判定和坐标输出，确认量化模型能够接入原有流程。

腕部纠偏则解决了全局定位与实际夹取之间的偏差。头部相机适合观察整盘目标，腕部相机适合接近目标后的局部测量。将连续对准移到 KV260 内，减少了图像回传和修正下发的等待；修正量由 PL 限幅、限速，再与轨迹求值结果叠加。大螺丝从悬停到收敛的用时 P50 为 1.06\,s，这一过程体现了板端视觉与硬件执行的配合。

实物调试还需要把停止运动与撤掉力矩分开处理。是否回到承托位置、关节反馈是否可信，会影响故障后的处置。因此，我们把异常检测、返回流程和人工介入一起设计，正常回位后结束使能，中途故障则按状态处理。保护条件和处置过程均保留运行记录，便于复查。

大模型工具参与了编码、测试、日志分析和文档整理。有效的协作方式是提供明确任务、接口和原始记录，对生成的修改做回归验证，并用实测检验结论。AMD Ross 的 HLS 技能帮助规范了可综合性检查和构建步骤。部署方法、协作规则及常见故障整理为附录中的技能包，供后续开发复用。
